% !TeX root = ../常微分方程笔记.tex
% 由主文档通过 \input 载入；请编译主文档。
\chapter{课程结构与学习路线}
\label{chap:ode-roadmap}

\section*{课程定位与编写范围}
本书面向已经学习一元微积分、正在学习数学分析和高等代数的读者，以本科常微分方程课程为主线，兼顾数学理论、求解方法和实际建模。研究对象是未知函数及其关于一个自变量的导数之间的关系。学习目标不仅是求出若干方程的解析表达式，还包括判断解是否存在、是否唯一、能够延伸到哪里，以及初值变化和时间演化如何影响解。

本书采用紧凑的教材写法，以定义、关键结论、必要证明和典型例题组织正文。前九章构成基础主线，后五章提供方法拓展与综合应用；标为“选学”的内容不作为基础章节的默认前提。篇幅较长的高阶理论只陈述所需版本并注明证明省略，其余核心推导尽量在正文中完成。

\section*{先修知识}
\begin{description}
  \item[微积分与数学分析] 导数与微分、定积分与微积分基本定理、泰勒公式、多元函数的偏导数、函数列的一致收敛。涉及压缩映射和级数解法时，再补充完备性、函数空间与幂级数。
  \item[线性代数] 向量空间、线性相关与线性无关、矩阵与行列式、线性方程组、特征值与特征向量。若尚未学习 Jordan 标准形，可先掌握可对角化情形，再补读重根情形。
  \item[辅助工具] 复数的指数表示、基本积分技巧与初步程序设计。程序设计仅用于数值实验，不作为理解基础理论的前提。所需知识集中列在附录\ref{app:ode-analysis}与附录\ref{app:ode-algebra}中。
\end{description}

\section*{全书课程大纲}
全书按照“提出问题与直接求解---初值问题的理论基础---线性结构---非线性定性分析---专题与综合应用”的顺序展开。
\begin{enumerate}
  \item \textbf{绪论与基本概念}（第\ref{chap:ode-intro}章）：建模、方程分类、解与解区间、初值问题、方向场。
  \item \textbf{一阶方程的初等解法}（第\ref{chap:ode-first-order}章）：变量分离、一阶线性方程、变量代换、恰当方程与积分因子。
  \item \textbf{隐式方程与可降阶方程}（第\ref{chap:ode-reduction}章）：导数未解出的方程、参数法、奇解与包络、降阶方法。
  \item \textbf{初值问题的存在性与唯一性}（第\ref{chap:ode-existence}章）：积分方程、Picard 迭代、压缩映射、存在唯一性与反例。
  \item \textbf{解的延拓与连续依赖性}（第\ref{chap:ode-dependence}章）：Gronwall 不等式、最大解、有限时间爆破、初值与参数依赖。
  \item \textbf{高阶线性微分方程}（第\ref{chap:ode-linear-scalar}章）：解空间、基本解组、常系数方程、非齐次方程、振动与共振。
  \item \textbf{线性微分方程组}（第\ref{chap:ode-linear-system}章）：基本矩阵、矩阵指数、Jordan 方法、常数变易与线性系统的长期行为。
  \item \textbf{自治系统与平面定性理论}（第\ref{chap:ode-phase-plane}章）：轨道、相图、平衡点、首次积分、周期轨道与极限环。
  \item \textbf{稳定性与 Lyapunov 方法}（第\ref{chap:ode-stability}章）：稳定性定义、线性化、Lyapunov 函数、不变性原理。
  \item \textbf{幂级数解法与特殊函数初步}（第\ref{chap:ode-series}章）：常点、正则奇点、Frobenius 方法、Bessel 与 Legendre 方程。
  \item \textbf{拉普拉斯变换与线性响应}（第\ref{chap:ode-laplace}章）：变换的存在条件、导数变换、卷积、分段输入与脉冲响应。
  \item \textbf{边值问题与 Sturm--Liouville 理论初步}（第\ref{chap:ode-boundary}章）：两点边值问题、Green 函数、特征值、正交性与振荡。
  \item \textbf{常微分方程数值方法}（第\ref{chap:ode-numerical}章）：Euler 与 Runge--Kutta 方法、误差、收敛、绝对稳定性与刚性。
  \item \textbf{数学建模与综合专题}（第\ref{chap:ode-modeling}章）：量纲与无量纲化、种群模型、振动模型、传染病模型及综合研究。
\end{enumerate}

\section*{学习顺序与内容层次}
第\ref{chap:ode-intro}至第\ref{chap:ode-stability}章构成基础主线，其中第\ref{chap:ode-reduction}章的参数法与奇解可先略读。第一次学习可先用第\ref{chap:ode-first-order}章积累求解经验，再学习第\ref{chap:ode-existence}、\ref{chap:ode-dependence}章的严格理论；第\ref{chap:ode-linear-scalar}、\ref{chap:ode-linear-system}章把微分方程与线性代数联系起来；最后由第\ref{chap:ode-phase-plane}、\ref{chap:ode-stability}章进入非线性问题。

第\ref{chap:ode-series}至第\ref{chap:ode-modeling}章是拓展与综合部分。级数解法需要幂级数知识；拉普拉斯变换主要承接线性方程；边值问题承接线性方程并引入新的边界约束；数值方法的算法部分可提前阅读，误差证明则依赖存在唯一性与连续依赖性。建模实例应贯穿全书，最后一章负责综合比较，不将应用完全留到最后。

\section*{统一记号与正文展开方式}
通常用 $t\in I\subseteq\mathbb{R}$ 表示自变量，其中 $I$ 是一个区间；用 $y(t)\in\mathbb{R}$ 表示标量未知函数，用 $x(t)\in\mathbb{R}^d$ 表示 $d$ 维状态向量。撇号表示对自变量求导，$y^{(n)}$ 表示 $n$ 阶导数，$\|\cdot\|$ 表示欧氏范数及相应的诱导矩阵范数。$t_0$ 表示初始时刻，$y_0$ 或 $x_0$ 表示初值。系数、参数和积分常数在各处另行声明。

正文围绕问题动机、概念与条件、关键推导和例题展开。习题覆盖概念辨析、计算、证明与反例及综合应用，附录给出部分答案与提示。学习时应区分“无法用已有方法写出显式解”与“没有解”，也应区分数值观察与已经证明的数学结论。

参考教材\cite{lebl-diffyqs,teschl-ode}可用于进一步阅读，并提供部分选学结论的完整证明。

